\section{Agent-facing interface checks}\label{sec:agentbench}

The agent-facing design described in Section~\ref{sec:agent} is the
most contestable claim in this paper: a reviewer may reasonably ask
whether a machine-readable schema \emph{actually} improves the
behaviour of an LLM agent on a real empirical task, or whether a
\emph{statsmodels-plus-prompting} baseline performs equivalently well.
That question needs frozen model releases, fixed prompts, matched tool
budgets, and a scoring rubric; it cannot be inferred from schema
coverage, and Section~\ref{sec:agentbench-behavioural} states what
answering it would take. The present section reports only the
\emph{mechanical} and \emph{contractual} evidence that the agent-facing
surface is correctly wired up. We refer to \citet{patil2023gorilla}
and \citet{liang2023helm} as methodological precedents for tool-use
and holistic-LLM benchmarks respectively.

\subsection{Mechanical evidence}\label{sec:agentbench-mechanical}

The function-schema layer is statically inspectable. We report in
Table~\ref{tab:schema-coverage} both the function-level coverage
($\RegistryTotal$ of $\RegistryTotal$ public functions emit a non-empty
OpenAI-style JSON Schema; \AgentCardCount{} currently expose curated or explicitly
inherited agent metadata cards) and the
parameter-level quality of those
schemas (\(\SchemaParameterTotal\) documented parameters across the
\(\RegistryTotal\) schemas, of which \(\SchemaParameterDefaultPct\) carry an
explicit default value and \(\SchemaParameterEnumPct\) carry enumerated
choices). Parameter quality
matters because the headline ``$\RegistryTotal/\RegistryTotal$ emit a schema'' figure is
high by construction -- the registry pipes \emph{any} typed signature
into a JSON Schema -- whereas the
stronger contract is that every parameter carries a
non-empty description. We do not turn
absence of package-distributed, package-wide OpenAI-style tool-schema catalogues
in \pkg{statsmodels}, \pkg{linearmodels}, or \pkg{scikit-learn} into
a numerical \emph{zero} score: those projects target different interfaces,
and Table~\ref{tab:schema-coverage} is an internal completeness audit,
not a behavioural or package-quality ranking; they do not rule out external
wrappers.

\begin{table}[t]
\centering\small
\begin{tabular}{@{}L{0.62\linewidth}rrr@{}}
\toprule
Audit item & Count & Missing & Share \\
\midrule
Public functions/classes                        & \RegistryTotal & --- & 100.0\% \\
Emit non-empty OpenAI-style tool schema          & \RegistryTotal & 0 & 100.0\% \\
Have at least one typed parameter                & \SchemaWithTypedParameter & --- & --- \\
Have non-empty function description              & \RegistryTotal & 0 & 100.0\% \\
Have curated or inherited agent metadata         & \AgentCardCount & \AgentDefaultCardCount & \AgentCardPct \\
\midrule
Parameters across all schemas                    & \SchemaParameterTotal & --- & 100.0\% \\
\quad with non-empty description                 & \SchemaParameterDescribed & 0 & \SchemaParameterDescribedPct \\
\quad with explicit default value                & \SchemaParameterDefault & --- & \SchemaParameterDefaultPct \\
\quad with enumerated choices (\code{enum})      & \SchemaParameterEnum & --- & \SchemaParameterEnumPct \\
\bottomrule
\end{tabular}
\caption{Internal mechanical schema audit for \statspai{} \StatsPAIVersion{}.
``Agent metadata'' means at least one non-empty assumption, pre-condition,
failure mode, limitation, or typical/minimum sample-size field after explicit
inheritance. The remaining schemas still emit name, signature, and
description fields, but are not counted as curated. Counts are regenerated
from the live registry; this table makes no cross-package quality claim.}
\label{tab:schema-coverage}
\end{table}

The same registry is exported as a versioned offline schema bundle
covering tool schemas, function schemas, agent cards, result-payload
validation, and provenance counts, so non-Python clients can discover
the tool surface without importing \statspai{}.

To keep this surface honest as the package grows, \statspai{} ships a
coverage ratchet for agent cards. The ratchet classifies registered
functions into cumulative documentation tiers, records per-field and
per-tier counters in a checked JSON floor, and fails CI on any
regression. The tiers deliberately separate a minimal schema-ready
description from richer assumptions, failure modes, alternatives, sample
size guidance, and certified/validated evidence. Lowering the floor
requires an explicit hand edit, so a schema-quality regression is
visible in code review rather than hidden behind a large function
count.

\subsection{Contractual trace evidence}\label{sec:agentbench-procedural}

The agent-facing surface is exposed through an MCP server
(Section~\ref{sec:mcp}) that negotiates Model Context Protocol
revisions 2024-11-05, 2025-03-26, and 2025-06-18. The protocol test
suite exercises the \code{initialize} handshake and version
negotiation, \code{tools/list} and \code{tools/call} with structured
content, and \code{isError} envelopes; the schema and registry suites
check that registered functions emit tool schemas, that agent-card
filters behave as documented, and that MCP tool discovery sees the same
registry-backed surface as \code{sp.function\_schema()}. These are the
package's own tests, not an external certification of protocol
conformance.

For this submission we also execute one deterministic end-to-end trace
through the live MCP server. The script
\code{replication/scripts/ex07\_agent\_trace.py} writes a temporary
\texttt{mpdta.csv} (the calibrated replica of Section~\ref{sec:ex-csdid},
so the traced estimate \(-0.0330\) is the replica value, not the
original-data \(-0.03995\)), calls \code{tools/list}, fits
\code{callaway\_santanna} with \code{as\_handle=true}, sends the
returned handle to \code{audit\_result}, and then sends the same
handle to \code{honest\_did\_from\_result}. It finally resolves the
returned citation keys through the MCP \code{bibtex} tool against the
curated \code{paper.bib}, so the trace exercises the citation
anti-hallucination contract rather than merely reporting key strings.
The same trace also sends a stale result-handle call to
\code{audit\_result} and checks that the MCP response is a recoverable
\code{isError} envelope with an actionable refit hint, so failure
recovery is exercised rather than only described.
Table~\ref{tab:agent-trace}
summarizes the trace; the normalized transcript is included in the
replication archive. The trace verifies the agent-specific claim that
the estimate \(\to\) audit \(\to\) sensitivity chain is executable
without ferrying arrays or hand-copying citations through the model
context.

\begin{table}[t]
\centering\small
\begin{tabular}{ll}
    \toprule
    Trace item & Value \\
    \midrule
    MCP tools listed & 559 \\
    Estimate & -0.0330 (SE 0.0078) \\
    Audit checks & 4 \\
    Honest-DiD max rejecting M & 0.0074 \\
    Recoverable stale-handle error & isError=True; hint present \\
    Verified citation keys & callaway2021difference, rambachan2023more \\
    BibTeX entries resolved & 2 \\
    \bottomrule
\end{tabular}

\caption{Deterministic MCP trace used for the agent-facing interface
check, including one recoverable stale-handle error path. Ephemeral
result handles are normalized in the transcript shipped with the
replication archive
(\code{replication/\allowbreak results/\allowbreak ex07\_agent\_trace.txt}).}
\label{tab:agent-trace}
\end{table}

\subsection{No behavioural benchmark claim}\label{sec:agentbench-behavioural}

Mechanical and contractual evidence does not answer whether an LLM agent
using this interface produces \emph{better empirical analyses} than one
using \pkg{statsmodels} or \proglang{R} packages through a Jupyter MCP
shim. That question requires frozen model
releases, fixed prompts, matched tool budgets, and a scoring rubric,
and belongs in a separate behavioural study. To make that deferral auditable
rather than vague, the archive includes
\code{replication/\allowbreak results/\allowbreak agent\_benchmark\_protocol.md}, which records
the matched arms, task families, scoring dimensions, leakage controls,
and minimum report items required before any behavioural claim is made.
The JSS contribution here is contractual: schemas, typed
errors, result handles, citations, and the deterministic trace are
reproducible interfaces, regenerated from the same validation ledger
that governs human-facing calls rather than from a separate
prompt-engineering layer, and they keep later benchmark scores from
being confounded with missing metadata or unverifiable citations.
